% Template: the frames for one chapter.
% Copy to documentation/tex/slides/sections/<topic>.tex, then add
%   \section{Chapter title}
%   \include{sections/<topic>}
% to slides/main.tex, in the same position the chapter has in the notes.

\begin{frame}
  \frametitle{Outline}
% One \nocite per line -- a multi-key \nocite spanning lines breaks bibtex.
\nocite{Bou18tra}
\nocite{CJP15alg}
Two or three sentences saying what this chapter is about,
in the same voice as the notes.

Say what the student should be able to do afterwards.
\end{frame}

\subsection{First idea}

\begin{frame}
  \frametitle{First idea}
One idea per frame.
Keep the lines short enough to read from the back of the room.

\pause

Use \pause between beats,
so that the frame is revealed in the order you talk through it.
\end{frame}

\begin{frame}
  \frametitle{A statement}
\begin{defi}
\label{def.something}
Reuse the label the notes give this statement,
so the two stay in step.
\[
x_t := 0.
\]
\end{defi}

\pause

Unnumbered display maths on slides uses \[ ... \],
not the equation environment.
\end{frame}

\begin{frame}
  \frametitle{What to take away}
Close the chapter with the two or three things
that should survive the lecture.
\end{frame}
